\subsection{What the literal manuscript does and does not say}

We use the numbering of the revised theorem on cote 131, scan 14, rather
than the different numbering on scan 2. Conditions 1--4 are Zariski locality,
fpqc descent, filtered-colimit compatibility on rings, and effectivity of
complete noetherian local rings from their Artinian quotients. Condition 5
is universal prorepresentability on Artinian bases with finite representing
local algebras. Condition 6 is valuative uniqueness. Conditions 7 and 8
concern local and generic spreading of modularity.

For an $S$-scheme $X$, we write $h_X(T)=\Hom_S(T,X)$ on the category of
$S$-schemes. A local model $(T,\eta)$ is \emph{locally modular} if it has
the universal property on local test schemes and the minimal
field-of-definition property of scan 35. When $F_{S_s}=h_{X_s}$, this is
equivalent to $\OO_{X_s,x}\to\OO_{T,t}$ being an isomorphism, as stated
on scan 37. Pro-modularity, which only tests Artinian neighbourhoods, is
a different property.

\begin{proposition}[The missing finiteness on the test scheme in 8]\label{A:8-literal}
\status{Refuted literally; necessity restored with locally finite type tests}
The necessity of condition 8 on scan 14 is false if ``$Z$ over $S'$''
allows arbitrary schemes, as its displayed wording does. If $Z\to S'$
is required to be locally of finite type, representability by a locally
quasi-finite separated scheme does imply this condition.
\end{proposition}
\begin{proof}
Let $k$ be an infinite algebraically closed field, $A=k[t]$, and
\[
 M=\bigoplus_{a\in k} A/(t-a),\qquad B=A\oplus M,
 \qquad (a,m)(b,n)=(ab,an+bm).
\]
Set $S'=S=\Spec A$, $Z=\Spec B$, and $F=h_S$. The square-zero ideal
$M$ gives a homeomorphism $Z\to S$. At the generic point its
localization is zero, so the structural morphism is locally modular
there: the local-ring map is $k(t)\xrightarrow{\sim}k(t)$.
At the point above $a\in k$, however, its local-ring map is
\[
 A_{(t-a)}\longrightarrow A_{(t-a)}\oplus k,
\]
which is not an isomorphism. Thus no nonempty open of $Z$ consists of
modular points. This violates 8 although $F$ is represented by the
identity $S\to S$, a finite separated morphism. The affine formula
defines the entire example; it uses an infinite module, so no finite
exhaustive computation is asserted.

For the corrected necessity, base change a representing scheme to
$S'$ and write $u:Z\to X_{S'}$ for the section. Both source and target
are locally of finite type over the noetherian $S'$, hence $u$ is locally
of finite presentation. If its local-ring map at $z$ is an isomorphism,
$u$ is etale there. The inverse local-ring isomorphism defines a
section over $\Spec\OO_{X_{S'},u(z)}$. Finite presentation spreads this
section to a neighbourhood of $u(z)$. After restricting to the etale
locus, its image is open, and $u$ identifies that image with an open
of $X_{S'}$. It is consequently locally modular on a nonempty open
of $Z$, as required.
\end{proof}

We have inspected scan 14 itself: it qualifies $S'$ as noetherian and
irreducible but adds no finite-type qualification to $Z$. The earlier
version on scan 2 does use a finite test scheme. This literal defect
does not settle the intended sufficiency theorem with finite-type tests;
that theorem, including its constructibility input, remains a separate
gate in Part~II.

\begin{proposition}[Embedded associated primes in NR(b)]\label{A:NR-counterexample}
\status{Refuted, with the correction in Proposition~\ref{A:NR-density}}
The literal NR(b) on cote 131, scan 16, is not a necessary condition for
representability by an unramified separated scheme. In particular,
``irreducible'' cannot replace schematic density of the generic point.
\end{proposition}
\begin{proof}
Let $k$ be any field and put
\[
 A=k[[t,\epsilon]]/(\epsilon^2,t\epsilon),\qquad
 \mathfrak p=I=(\epsilon),\qquad S=\Spec A,\qquad X=\Spec(A/I).
\]
The ring $A$ is complete noetherian local with maximal ideal
$(t,\epsilon)$. Its nilradical is $\mathfrak p$, since $A/\mathfrak p=k[[t]]$
is a domain. Thus $\Spec A$ is irreducible of dimension one and
$I\mathfrak p=0$. The $A$-module description
$A=k[[t]]\oplus k\epsilon$ proves that $\epsilon\ne0$.
The morphism $X\to S$ is a finitely presented closed immersion, hence
finite, unramified, and separated.

Set $F=h_X$. The identity of $X$ defines $\eta\in F(A/I)$.
In $A_{\mathfrak p}$ the element $t$ is invertible, so $\epsilon=0$ and
$A_{\mathfrak p}=k((t))=(A/I)_{\mathfrak p}$. Therefore the restriction
of $\eta$ has a lift in $F(A_{\mathfrak p})$. But $F(A)$ is empty:
an $A$-algebra map $A/I\to A$ would have to send the image of $\epsilon$,
which is zero, to the nonzero structural element $\epsilon$.
This verifies every displayed hypothesis of NR(b) and contradicts its conclusion.
The annihilator of $\epsilon$ is $(t,\epsilon)$, displaying the embedded
associated prime which causes the failure.
\end{proof}

This example was already proposed, conditionally, in the supplied
\texttt{MPL\_131\_notes.md}, RN020; we have re-proved it here and inspected
scan 16. The preceding RN019 correctly distinguished topological
irreducibility from schematic density. The later study's description of the
test rings as ``domains'' changes the statement: in a domain $\mathfrak p=0$
and its displayed test has no nonzero $I$.

\begin{proposition}[The density correction]\label{A:NR-density}
\status{Proved under stated hypotheses}
Let $(A,\mathfrak m)$ be a complete noetherian local ring of dimension one
with exactly one associated prime $\mathfrak p$, and let $I\subseteq
\mathfrak p$ satisfy $I\mathfrak p=0$. Let $X\to\Spec A$ be an
unramified separated morphism of schemes. An $A$-morphism
$\eta_0:\Spec(A/I)\to X$ whose generic restriction extends to
$\Spec A_{\mathfrak p}\to X$ extends uniquely to $\Spec A\to X$.
\end{proposition}
\begin{proof}
Because $\mathfrak p$ is the sole associated prime, every element outside
$\mathfrak p$ is a non-zero-divisor. Hence $A\to A_{\mathfrak p}$ is
injective. Write $x$ for the image of the closed point under $\eta_0$.
An unramified morphism factors, in a neighbourhood of $x$, as a closed
immersion into a scheme $E$ etale over $A$ (Stacks, Proposition 10.152.1,
Tag 0395, applied to affine neighbourhoods) \cite{Stacks0395}. The inverse image of that
neighbourhood in the local scheme $\Spec(A/I)$ contains its closed point
and is therefore the whole scheme. Since $I^2=0$, formal etaleness gives
a unique lift $\widetilde\eta:\Spec A\to E$ of the composite of $\eta_0$
with $X\hookrightarrow E$.

Let $J$ be the ideal of the pullback $X\times_E\Spec A\subseteq\Spec A$.
It lies in $I$. On $\Spec A_{\mathfrak p}$ the given lift to $X$ and
$\widetilde\eta$ have the same reduction modulo $I_{\mathfrak p}$;
formal unramifiedness of $E$ makes their composites to $E$ equal.
Thus $J A_{\mathfrak p}=0$. Injectivity of $A\to A_{\mathfrak p}$
forces $J=0$, so $\widetilde\eta$ factors through $X$. Uniqueness follows
from formal unramifiedness of $X$. The factorization and the argument may
also be made after an etale neighbourhood; the residue-field equality
provided by the section and the henselian property of $A$ provide the
required local etale branch.
\end{proof}

Murre's published condition (F7), in Expose 294, pp.~245--246,
uses precisely a unique associated prime \cite{Murre1965}.
The earlier wording in cote 131, scan 7, explicitly requires schematic
density. Proposition~\ref{A:NR-density} explains the exact role of that
hypothesis rather than inferring it from irreducibility.

\subsection{The Artinian test extension and the etale corollary}

\begin{lemma}[Extending infinitesimal tests]\label{A:artinian-tests}
\status{Proved}
Let $S$ be locally noetherian and let $F$ satisfy conditions 1--4 of
scan 14. If $F(A)\to F(A/I)$ is injective for every Artinian local
$S$-algebra $A$ and nilpotent ideal $I$, it is injective for every
$S$-algebra $A$ and nilpotent ideal $I$. If all the Artinian maps are
bijective, all these maps are bijective.
\end{lemma}
\begin{proof}
We work over an affine open $\Spec R$ of $S$; Zariski locality reduces
the general case to this one. First take a noetherian $R$-algebra $A$
and two elements whose reductions modulo $I$ agree. For each maximal
ideal $\mathfrak m$ of $A$, their images in
$F(A_{\mathfrak m}/\mathfrak m^n)$ agree by the Artinian injectivity
hypothesis. Condition 4 makes their images in
$F(\widehat{A_{\mathfrak m}})$ equal. Completion is faithfully flat,
so condition 2 makes their images in $F(A_{\mathfrak m})$ equal.
By condition 3 this equality holds on some principal neighbourhood of
$\mathfrak m$ in $\Spec A$. Such neighbourhoods cover $\Spec A$:
every prime specializes to a maximal ideal, and opens are stable under
generization. Condition 1 gives equality in $F(A)$.

For an arbitrary $R$-algebra $A$, write it as the filtered union of
its finite-type $R$-subalgebras $A_i$, and put $I_i=A_i\cap I$.
These rings are noetherian, $I_i$ is nilpotent, and
$A/I=\varinjlim_i A_i/I_i$. Condition 3 descends the two elements and
the equality of their reductions to a common finite stage. The result
just proved there gives injectivity for $A$. This step is established
before we use any nonnoetherian tensor-product ring.

Now assume the Artinian maps are bijective. Let $A$ be noetherian local
and $\eta_0\in F(A/I)$. The images of $\eta_0$ in
$F(A/(I+\mathfrak m^n))$ lift uniquely to $F(A/\mathfrak m^n)$.
Uniqueness makes these lifts compatible, and condition 4 supplies
$\widehat\eta\in F(\widehat A)$ whose reduction is the pullback of
$\eta_0$. Here $\widehat A/I\widehat A$ is the complete quotient,
so condition 4 also verifies this last equality. The two pullbacks of
$\widehat\eta$ to $F(\widehat A\otimes_A\widehat A)$ agree modulo
the nilpotent ideal generated by $I$. The injectivity already proved
for \emph{all} rings therefore makes them equal. Condition 2 descends
$\widehat\eta$ to $\eta\in F(A)$, and faithful flatness checks its
reduction.

For a general noetherian $A$, apply this to all $A_{\mathfrak m}$.
Condition 3 spreads each lift and its prescribed reduction to a
principal neighbourhood; the injectivity already proved makes the
lifts agree on overlaps, and condition 1 glues them. Finally the
same pairs $(A_i,I_i)$ as above descend any $\eta_0\in F(A/I)$ to
a noetherian stage. Lifting there and pulling back proves existence
over arbitrary $A$; uniqueness is the first half of the proof.
\end{proof}

\begin{theorem}[Corollary 1 of cote 131]\label{A:etale-corollary}
\status{Known, with the Artinian reduction proved here}
For $S$ locally noetherian, a functor $F:(\mathrm{Sch}/S)^{\mathrm{op}}
\to\mathrm{Sets}$ is represented by an etale separated $S$-scheme if
and only if it satisfies conditions 1, 2, 3, 4, 6 of scan 14 and is
formally etale on Artinian local test algebras. Thus the Artinian
alternative explicitly written in Corollary 1 on scan 16 is sufficient.
\end{theorem}
\begin{proof}
For necessity, an etale morphism is formally etale; its functor has
Zariski and fpqc descent, commutes with filtered ring colimits and
complete noetherian local adic limits, and satisfies valuative
uniqueness when separated. For sufficiency,
Lemma~\ref{A:artinian-tests} turns the Artinian hypothesis into formal
etaleness for every nilpotent ideal in every ring. Murre, Expose 294,
Corollary 2, p.~247, applies with his (F1) equal to our 1 and 2,
(F2) to 3, (F3) to 4, (F6) to 6, and his strengthened (F5) equal
to this formal etaleness \cite{Murre1965}. Its conclusion is a
locally finite-type, etale, separated scheme. The base is the same
locally noetherian $S$ throughout; no excellence or finite-type
hypothesis on $S$ has been introduced.
\end{proof}

\subsection{A dimension-one example making condition 8 indispensable}

\begin{theorem}[Finite local models without generic spreading]\label{A:ind-thickening}
\status{Proved}
Let $k$ be an infinite algebraically closed field and $S=\Spec k[t]$.
For each finite subset $E\subset k$, put
\[
 f_E(t)=\prod_{a\in E}(t-a),\qquad
 X_E=\Spec k[t,x]/(x^2,f_E(t)x).
\]
If $E\subseteq E'$, the evident quotient of coordinate rings defines
a closed immersion $X_E\hookrightarrow X_{E'}$. Let $F$ be the
Zariski sheaf associated to $\varinjlim_E h_{X_E}$. Then $F$ satisfies
conditions 1--7 of the revised theorem, is represented over every local
base scheme by a finite separated scheme, but is not represented by a
scheme locally of finite type over $S$. The zero section
$\eta:S\to F$ is locally modular exactly at the generic point of $S$.
Consequently it fails condition 8, even with the finite test $Z=S$.
\end{theorem}
\begin{proof}
For an affine $S$-scheme $T=\Spec B$, with structural element $t_B\in B$,
the sheaf has the explicit formula
\[
 F(B)=\{b\in B:\ b^2=0\text{ and }f_E(t_B)b=0
                 \text{ for some finite }E\subset k\}.
\]
Indeed, finitely many local witnesses on an affine scheme can be replaced
by their union $E$; their polynomial relations then hold globally.
This formula realizes $F$ as a subfunctor of $h_{\mathbb A^1_S}$.
It proves filtered-colimit compatibility: an element and its two equations
descend to one finite stage, and equality is detected at a finite stage.
It also proves fpqc descent. If $B\to B'$ is faithfully flat and an element
$b'\in F(B')$ has equal pullbacks to $B'\otimes_B B'$, affine descent
gives $b\in B$; a finite witness $E$ for $b'$ works for $b$ by faithful
flatness. Combined with Zariski locality this is the fpqc sheaf property.
Valuative uniqueness follows from the injection $V\hookrightarrow
\operatorname{Frac}(V)$ for every valuation domain $V$.

Now let $B$ be any local $k[t]$-algebra, with residue field $L$.
There is at most one $a\in k$ with $t_B-a$ noninvertible in $B$.
If there is none, every $f_E(t_B)$ is a unit; all the base changes
$(X_E)_B$ equal $\Spec B$. If there is such an $a$, the cofinal subsystem
of subsets $E$ containing $a$ becomes the constant scheme
\[
 Y_B=\Spec B[x]/(x^2,(t_B-a)x).
\]
Here the other factors of $f_E(t_B)$ are units even after arbitrary
base change from $B$. Thus these are isomorphisms of functors on all
$B$-schemes, not merely equalities on $B$-valued points. Each $Y_B$ is
finite and separated over $B$.

In particular, the Artinian prorepresentability and finite local-algebra
condition holds. If $B$ is complete noetherian local, the same finite
model represents the functor on $B$ and all $B/\mathfrak m^n$; its
coordinate algebra is finitely presented, so maps into it are exactly
compatible solutions of its polynomial equations at all these levels.
This proves condition 4 with its actual comparison map.

For condition 7, base change to a complete noetherian local $A$.
The functor is represented by a finite $A$-scheme $Y_A$. If
$\Spec B\to Y_A$ is locally modular at the closed point, where
$B$ is a finite local $A$-algebra, its map to the local component of
$Y_A$ induces an isomorphism on local rings. A finite scheme over a
complete local ring is a disjoint union of local schemes, so this map
identifies $\Spec B$ with an open-and-closed component. It is therefore
locally modular at every point. This verifies condition 7.

At the generic point, the zero section identifies both local models with
$\Spec k(t)$ and is locally modular. At a closed point $a$, its map on
local rings is
\[
 k[t]_{(t-a)}[x]/(x^2,(t-a)x)\longrightarrow k[t]_{(t-a)},\qquad x\mapsto0.
\]
The class of $x$ is nonzero, so this is not an isomorphism. Thus the
modularity locus is the singleton generic point, which contains no
nonempty open of $S$.

Finally suppose $F=h_X$ for a scheme $X$ locally of finite type over $S$.
The zero section would be an $S$-morphism $s:S\to X$. Around its generic
point, choose an affine finite-type neighbourhood $W\subseteq X$ and
shrink the domain to an affine nonempty open $U\subseteq S$ mapping
into $W$. The coordinate map $\OO(W)\to\OO(U)$ is surjective, since it
splits the structural map from $\OO(U)$ after replacing $W$ by its
base change to $U$. Its kernel is finitely generated and vanishes after
tensoring with $k(t)$, since $X_{k(t)}=\Spec k(t)$. After one further
shrinking of $U$, that kernel is zero. Hence the section is an
isomorphism onto an open of $X$ over $U$, contradicting the local
calculation at every closed point of $U$. This also proves
nonrepresentability without assuming the openness conclusion that
condition 8 is intended to supply.
\end{proof}

\subsection{Completion descent and the nonfinite quasi-projective question}

\begin{theorem}[The descent step actually needed in cote 131]\label{A:completion-descent}
\status{Known, with hypotheses checked}
Let $(A,\mathfrak m)$ be noetherian local and let $\widehat A$ be its
$\mathfrak m$-adic completion. Every descent datum on a separated,
locally quasi-finite $\widehat A$-scheme $X'$ relative to
$\Spec\widehat A\to\Spec A$ is effective as a scheme. The descended
scheme $X$ is separated and locally quasi-finite over $A$.
\end{theorem}
\begin{proof}
Murre, Expose 294, Section 4, Theorem, case (d), proves effectivity for a
faithfully flat quasi-compact $S'\to S$ with $S'$ locally noetherian and
$X'\to S'$ separated, locally quasi-finite, and locally of finite
presentation \cite{Murre1965}. Completion of a noetherian local ring is
noetherian and faithfully flat; the induced affine morphism is
quasi-compact. Over the noetherian base $\widehat A$, local finite type
implies local finite presentation. Thus all its assumptions hold here.
The comparison isomorphism is
$X\times_{\Spec A}\Spec\widehat A\simeq X'$, compatible with the given
cocycle. Separation and local quasi-finiteness descend under faithful
flat quasi-compact base change; for the latter, see
Stacks, Lemma 35.23.26, Tag 02VI \cite{Stacks02VI}.
This proves lead L3 in the generality required by the manuscript's step 3.
\end{proof}

For quasi-compact pieces, the route proposed in L3 is also valid: a
quasi-finite separated scheme over an affine base is quasi-affine, and
quasi-affine schemes satisfy effective fpqc descent. Passing from arbitrary
pieces to a covering by \emph{descent-invariant} pieces requires an argument;
quasi-compactness of an arbitrarily selected open is insufficient. The
opened theorem above supplies that argument under the noetherian
hypotheses actually used. Merely saying ``descend as an algebraic space''
would not by itself prove effectivity along the completion map.

\begin{theorem}[The nonfinite quasi-projective descent question]\label{A:nonfinite-descent}
\status{Refuted, with the replacements below}
In cote 126, scans 32--35, deleting all but one point of a fibre of a
finite etale cover need not leave a morphism of effective descent for
quasi-projective schemes. The failure already occurs for a projective
scheme over a surjective, quasi-compact etale cover of a smooth complex
scheme.
\end{theorem}
\begin{proof}
We use the construction of SGA~1, VIII, 7.10(ii), with explicit equations
from Stacks, Lemmas 110.67.1 and 110.68.1
\cite{SGA1VIII7,Stacks08J1,Stacks08KF}. Put
\[
 g[x,y,z,w]=[y,x,w,z],\quad
 Y=\mathbf P^3_{\mathbf C}\setminus
 \bigl(\{[x,x,z,z]\}\cup\{[x,-x,z,-z]\}\bigr).
\]
The involution $g$ acts freely, so $Y\to S=Y/\langle g\rangle$ is a
finite etale double cover of a smooth quasi-projective scheme.
The conics
\[
 C=V(xy-z^2,w),\qquad D=V(xy-w^2,z)
\]
lie in $Y$, are exchanged by $g$, and meet transversely at
$P=[1,0,0,0]$ and $Q=gP$. Over $Y-Q$ blow up $C$ and then the strict
transform of $D$; over $Y-P$ reverse their order. Away from $P,Q$ the
centres are disjoint, so these constructions glue canonically to a
smooth scheme $V_Y$, proper over $Y$. The involution lifts and supplies
a descent datum.

The quotient $V_Y/\langle g\rangle$ is not a scheme. Here is the
geometric obstruction used in the cited construction. In the proper
exceptional surface above $C\cup D$, write the fibre over $P$ as
$L_0\cup M_0$ and the fibre over $Q$ as $L'_0\cup M'_0$, with
$gL_0=M'_0$. If $L,M$ are general fibres over $C,D$, the two blowup
orders give
\[
 L\sim L_0+M_0\sim L'_0,\qquad
 M\sim M_0\sim L'_0+M'_0.
\]
Thus $L_0+gL_0\sim0$. If the quotient were a scheme, the image of
these two disjoint curves would be a projective line. At a point of
that line, a local function vanishing at the point but not identically
on the line gives, by taking the closure of a codimension-one component,
an effective divisor meeting but not containing the line. Its pullback
to the smooth $V_Y$ is Cartier and has positive degree on
$L_0+gL_0$, contradicting the displayed rational equivalence.

Now set $T'=Y-Q$. It surjects onto $S$: it contains $P$ above the
exceptional image $s$, and contains both points above every other
closed orbit. It is etale and quasi-compact. It is not finite, since
its fibre has size one at $s$ and size two generically, whereas a finite
etale morphism over the connected $S$ has constant rank. The restriction
$V'=V_Y|_{T'}$ is projective, being an iterated blowup. Its datum on
$T'\times_S T'$ is the restriction of the lifted involution; the
off-diagonal part is defined on $Y-\{P,Q\}$. If this datum were
effective as a scheme, faithfully flat descent of morphisms would
identify its base change to $Y$ with $V_Y$, hence its result with the
impossible scheme quotient above. This is exactly the operation of
retaining one point in a fibre asked for in the manuscript.
\end{proof}

\begin{proposition}[Effective replacements]\label{A:polarized-descent}
\status{Known, with hypotheses checked}
Let $S'\to S$ be faithfully flat and quasi-compact.
\begin{enumerate}
\item If it is finite locally free, every descent datum on a
quasi-projective $S'$-scheme is effective, and its result is
quasi-projective over $S$.
\item For arbitrary $S'\to S$, a descent datum on a quasi-compact
$S'$-scheme $X'$ together with a compatible datum on an $S'$-ample
invertible sheaf $\mathscr L'$ is effective. The descended invertible
sheaf is $S$-ample.
\end{enumerate}
\end{proposition}
\begin{proof}
These are SGA~1, VIII, Corollary 7.7 and Proposition 7.8
\cite{SGA1VIII7}. In (1), norming an ample bundle over the finite
locally free cover gives invariant affine neighbourhoods and a
descended ample bundle. In (2), the compatible datum descends the
graded algebra of sections of powers of $\mathscr L'$. Relative
$\operatorname{Proj}$ commutes with flat base change; the canonical
open immersion of $X'$ into that relative $\operatorname{Proj}$
is invariant and descends. The inverse-image comparison is exactly
$X\times_S S'\simeq X'$, with the prescribed cocycle. These statements
apply to the two corrected hypotheses as written; the Hironaka
example has neither finiteness of the cover nor a compatible ample
bundle.
\end{proof}
